\clearpage
\section*{补充推导：固定事件下的热能耗表达}
\label{sec:formal-details}

为说明第3.5节条件传播与积分系数的来源，下面在给定联合工作点和固定片长下展开对应关系。正文的调度选项和频率控制约束负责确定片段的实际状态；本部分推导不新增任务或运行假设。

\subsection*{一、固定控制状态下的物理系数}

在联合工作点状态$\omega$下，各功率组的基准及温度漏电系数由该状态决定。参考温度偏置与节点基准注热向量可写为
\begin{equation}
(b^P_\omega)_n=
\sum_g a^B_{ng}(\xi_g(\omega))B_g(\xi_g(\omega))
+\sum_g\lambda_{gn}(\xi_g(\omega))(T_a-T_g^{\rm ref}).
\label{eq:baseline-vector}
\end{equation}
记$(Q_c)_{no}=h_{n,u_o}q_o$，并令$\ell_{\omega,n}=\sum_g\lambda_{gn}(\xi_g(\omega))$，则节点注热、状态方程系数和片内电功率可整理为
\begin{equation}
\begin{aligned}
p^{\rm heat}&=b^P_\omega+\operatorname{diag}(\ell_\omega)\theta+Q_cv,\\
b_\omega&=C^{-1}b^P_\omega,\qquad D_c=C^{-1}Q_c,\\
\beta_\omega&=\mathbf1^{\mathsf T}b^P_\omega,\\
P^{\rm net}&=\beta_\omega+\ell_\omega^{\mathsf T}\theta+q_c^{\mathsf T}v.
\end{aligned}
\label{eq:physical-coefficients}
\end{equation}
其中$\mathbf1^{\mathsf T}Q_c=q_c^{\mathsf T}$。$A_\omega$的单位为s$^{-1}$，$b_\omega$及$D_c$给出温度变化率，$\beta_\omega$、$q_c$的单位为W，$\ell_\omega$的单位为W/K。上述关系同时对应式\eqref{eq:local-heat-input}的注热和式\eqref{eq:group-electric-power}的计能，防止一项物理功率在两处重复计算。

\subsection*{二、固定片段内的温度传播与能量积分}

当联合状态$\omega$及片段活动向量$v$均不变时，热方程为常系数仿射系统。对长度为$\delta$的片段，定义
\begin{equation}
\begin{aligned}
\Phi_\omega(\delta)&=e^{A_\omega\delta},\\
\mathsf K_\omega(\delta)&=\int_0^\delta e^{A_\omega t}\,dt,\\
\mathsf L_\omega(\delta)&=\int_0^\delta(\delta-t)e^{A_\omega t}\,dt.
\end{aligned}
\label{eq:matrix-integrals}
\end{equation}
温度轨迹因此为
\begin{equation}
\theta(t)=e^{A_\omega t}\theta(0)
+\mathsf K_\omega(t)(b_\omega+D_cv).
\label{eq:segment-solution}
\end{equation}
将温度解代入片内电功率并积分，可得
\begin{equation}
\begin{aligned}
\theta(\delta)&=\Phi_\omega(\delta)\theta(0)
+\mathsf K_\omega(\delta)b_\omega
+\mathsf K_\omega(\delta)D_cv,\\
\Delta E&=\ell_\omega^{\mathsf T}\mathsf K_\omega(\delta)\theta(0)
+\beta_\omega\delta+\ell_\omega^{\mathsf T}\mathsf L_\omega(\delta)b_\omega\\
&\quad+\left[\delta q_c^{\mathsf T}
+\ell_\omega^{\mathsf T}\mathsf L_\omega(\delta)D_c\right]v.
\end{aligned}
\label{eq:segment-integrals}
\end{equation}
与正文式\eqref{eq:thermal-transition}及式\eqref{eq:energy-transition}对照，传播与积分系数分别为
\begin{equation}
\begin{aligned}
\psi&=\mathsf Kb_\omega,&\Gamma&=\mathsf KD_c,\\
\mu^{\mathsf T}&=\ell_\omega^{\mathsf T}\mathsf K,&
\zeta&=\beta_\omega\delta+\ell_\omega^{\mathsf T}\mathsf Lb_\omega,\\
\nu^{\mathsf T}&=\delta q_c^{\mathsf T}+\ell_\omega^{\mathsf T}\mathsf LD_c.
\end{aligned}
\label{eq:integral-coefficients}
\end{equation}
所有系数均在当前$\omega$、$\delta_{c\ell}$下取值。该表达不需要计算$A_\omega^{-1}$，因而不要求该矩阵可逆。阶段活动、等待与最低档空闲尾段都遵循同一传播公式，只是活动向量和联合工作点有所不同。

\subsection*{三、有效变量界与条件线性化}

有限事件、有限工作点和固定窗口保证相应系数存在有限界。采用相容无穷范数，定义
\begin{equation}
a_* =\max_{\omega\in\Omega}\|A_\omega\|_\infty,
\qquad b_{c*}=\max_{\omega\in\Omega}
\left(\|b_\omega\|_\infty+\|D_c\|_\infty\right).
\label{eq:norm-coefficients}
\end{equation}
因为$\|v\|_\infty\le1$，积分比较可给出对整个窗口有效的保守界
\begin{equation}
\|\theta_c(t)\|_\infty\le H_c,
\quad H_c=e^{a_*W_c}
\left(\|T_c^0-T_a\mathbf1\|_\infty+b_{c*}W_c\right).
\label{eq:state-bound}
\end{equation}
该界对任意合法模式和状态切换成立，包括尚未满足温限的时间表。令
\begin{equation}
\overline P_c=\max_\omega
\left(|\beta_\omega|+\|\ell_\omega\|_1H_c+\|q_c\|_1\right),
\quad|\Delta E_{c\ell}|\le\delta_{c\ell}\overline P_c.
\label{eq:energy-bound}
\end{equation}
则温度状态与段能量的有限变量盒可用于构造条件等式的残差界。例如，式\eqref{eq:thermal-transition}与式\eqref{eq:energy-transition}分别可取
\begin{equation}
\begin{aligned}
M^T_n&=H_c+\|\Phi_{n:}\|_1H_c+|\psi_n|+\|\Gamma_{n:}\|_1,\\
M^E&=\delta_{c\ell}\overline P_c+\|\mu\|_1H_c+|\zeta|+\|\nu\|_1.
\end{aligned}
\label{eq:residual-bounds}
\end{equation}
当状态指示变量为1，等式必须成立；为0时，残差界须覆盖所有允许的变量值。式中的大$M$用于精确表达条件关系，不是允许温度超过上限的容差。实际数值建模可收紧这些界，但不能无依据地排除合法的有限域安排。

对于工作点固定的片段，温度有效范围、组功率和节点注热的物理有效性可以统一写成
\begin{equation}
a_{j\omega}^{\mathsf T}\theta+b_{jc\omega}^{\mathsf T}v+d_{j\omega}\le0.
\label{eq:validity-affine}
\end{equation}
在片段起点和终点可以施加对应的必要约束，片段内部仍需连续检查。窗口末点没有正在执行的Stage，因此各控制域均为$L_h$，其功率有效性应按\emph{最低档联合状态}及$v=0$核查，不能沿用旧的“最后一次工作点”或“域状态保持”假设。对任意条件不等式$w=1\Rightarrow a^{\mathsf T}\theta+b^{\mathsf T}v+d\le0$，取$M=\|a\|_1H_c+\|b\|_1+|d|$即可写成对应的有界线性形式。若内点发现明确违规，可用同一状态下的温度解析式构建必要反馈行。

\subsection*{四、选项、状态与真实时间表之间的对应}

给定受限问题$P_{\mathcal G}$的一份合法安排，每个Stage实例都对应一个唯一模式--起点选项，因而能够构造满足任务分配、期限和资源互斥的$x,y$。真实域状态产生$\rho,w$：正在执行的Stage锁定其工作点，没有正在执行的Stage时非最低档支持式\eqref{eq:idle-auto-low}迫使最低档$L_h$。按真实温度轨迹填写片段端点温度和完整段能量，则各条件传播、积分等式均成立。

反过来，若一个整数解满足模式选择、依赖、互斥和域控制约束，便可恢复一张合法时间表。式\eqref{eq:idle-auto-low}与工作阶段锁定式共同确定每个片段的实际控制状态，既不允许空闲片段任意换至高档，也不允许有工作执行期间换档。给定初始温度后，常系数热方程在每段有唯一解，片段端点状态及能量恰由前述积分表达给出。再经完整连续温度与有效域检查合格，便得到$P_{\mathcal G}$的一份可行方案。

反馈行需要保留所有真实安全安排：若候选在某片段内部确有超温，则该时刻的温限是安全安排必须满足的必要条件；未采用对应联合状态的安排不受该行限制。完整选项赋值排除只应针对已经可靠确认违规的完整场景候选。数值未决不能被当作超温或不可行证明。上述结论是模型的数学对应关系，不自动赋予有限计算预算内的求解器结果全局最优证明。
